\section{Reproducing the results}\label{app:repro}
\begin{verbatim}
make                              # the five solvers, in build/
make test                         # heap tests, amortized check
python3 scripts/fuzz.py           # 300 random graphs, all solvers
python3 scripts/gen_instances.py  # generated instances
bash scripts/run_bench.sh         # results/*.csv
python3 scripts/make_tables.py    # report/generated/*
make report                       # report/report.pdf
\end{verbatim}
Every solver reads a SteinLib \texttt{.stp} file and prints one CSV row; \texttt{--tree=FILE} writes the tree's edges.
The README in the archive lists all options.

\section{Per-instance results}\label{app:tables}
Costs in bold equal the optimum. A superscript t marks a branch-and-bound run stopped by the 300\,s limit; its cost is the best tree
found. A dash in the DP columns means \(k>20\), which the DP does not attempt. Times are wall-clock seconds for one run.
\subsection*{Set B}
\begingroup\scriptsize\setlength{\tabcolsep}{3.2pt}
\begin{longtable}{@{}lrrrrrrrrrrrrrr@{}}\toprule
Inst. & \(n\) & \(m\) & \(k\) & OPT & \multicolumn{2}{c}{G1} & \multicolumn{2}{c}{G2} & \multicolumn{2}{c}{G3} & \multicolumn{2}{c}{DP} & \multicolumn{2}{c}{B\&B}\\
 & & & & & cost & s & cost & s & cost & s & cost & s & cost & s\\\midrule\endhead
b01 & 50 & 63 & 9 & 82 & \textbf{82} & 0.000 & \textbf{82} & 0.000 & \textbf{82} & 0.000 & \textbf{82} & 0.003 & \textbf{82} & 0.000\\
b02 & 50 & 63 & 13 & 83 & 88 & 0.000 & \textbf{83} & 0.000 & \textbf{83} & 0.000 & \textbf{83} & 0.100 & \textbf{83} & 0.000\\
b03 & 50 & 63 & 25 & 138 & \textbf{138} & 0.000 & \textbf{138} & 0.002 & \textbf{138} & 0.000 & -- & -- & \textbf{138} & 0.001\\
b04 & 50 & 100 & 9 & 59 & 62 & 0.000 & \textbf{59} & 0.002 & \textbf{59} & 0.000 & \textbf{59} & 0.002 & \textbf{59} & 0.003\\
b05 & 50 & 100 & 13 & 61 & \textbf{61} & 0.000 & \textbf{61} & 0.001 & \textbf{61} & 0.000 & \textbf{61} & 0.128 & \textbf{61} & 0.002\\
b06 & 50 & 100 & 25 & 122 & 127 & 0.002 & 124 & 0.006 & 124 & 0.000 & -- & -- & \textbf{122} & 0.010\\
b07 & 75 & 94 & 13 & 111 & \textbf{111} & 0.000 & \textbf{111} & 0.003 & \textbf{111} & 0.004 & \textbf{111} & 0.213 & \textbf{111} & 0.005\\
b08 & 75 & 94 & 19 & 104 & \textbf{104} & 0.000 & \textbf{104} & 0.001 & \textbf{104} & 0.000 & \textbf{104} & 27.2 & \textbf{104} & 0.001\\
b09 & 75 & 94 & 38 & 220 & 221 & 0.000 & \textbf{220} & 0.021 & \textbf{220} & 0.001 & -- & -- & \textbf{220} & 0.007\\
b10 & 75 & 150 & 13 & 86 & 91 & 0.000 & \textbf{86} & 0.001 & 87 & 0.000 & \textbf{86} & 0.145 & \textbf{86} & 0.003\\
b11 & 75 & 150 & 19 & 88 & 92 & 0.000 & \textbf{88} & 0.002 & \textbf{88} & 0.001 & \textbf{88} & 26.6 & \textbf{88} & 0.006\\
b12 & 75 & 150 & 38 & 174 & \textbf{174} & 0.002 & \textbf{174} & 0.012 & \textbf{174} & 0.003 & -- & -- & \textbf{174} & 0.049\\
b13 & 100 & 125 & 17 & 165 & 175 & 0.001 & \textbf{165} & 0.004 & 170 & 0.000 & \textbf{165} & 5.225 & \textbf{165} & 0.002\\
b14 & 100 & 125 & 25 & 235 & 237 & 0.005 & \textbf{235} & 0.017 & 236 & 0.007 & -- & -- & \textbf{235} & 0.003\\
b15 & 100 & 125 & 50 & 318 & \textbf{318} & 0.001 & \textbf{318} & 0.048 & \textbf{318} & 0.001 & -- & -- & \textbf{318} & 0.012\\
b16 & 100 & 200 & 17 & 127 & 135 & 0.000 & \textbf{127} & 0.007 & 133 & 0.001 & \textbf{127} & 4.922 & \textbf{127} & 0.024\\
b17 & 100 & 200 & 25 & 131 & 134 & 0.004 & 132 & 0.017 & 132 & 0.004 & -- & -- & \textbf{131} & 0.020\\
b18 & 100 & 200 & 50 & 218 & 221 & 0.002 & 220 & 0.061 & \textbf{218} & 0.002 & -- & -- & \textbf{218} & 0.037\\
\bottomrule\end{longtable}\endgroup
\subsection*{Set C}
\begingroup\scriptsize\setlength{\tabcolsep}{3.2pt}
\begin{longtable}{@{}lrrrrrrrrrrrrrr@{}}\toprule
Inst. & \(n\) & \(m\) & \(k\) & OPT & \multicolumn{2}{c}{G1} & \multicolumn{2}{c}{G2} & \multicolumn{2}{c}{G3} & \multicolumn{2}{c}{DP} & \multicolumn{2}{c}{B\&B}\\
 & & & & & cost & s & cost & s & cost & s & cost & s & cost & s\\\midrule\endhead
c01 & 500 & 625 & 5 & 85 & 88 & 0.001 & \textbf{85} & 0.002 & 87 & 0.002 & \textbf{85} & 0.010 & \textbf{85} & 0.030\\
c02 & 500 & 625 & 10 & 144 & \textbf{144} & 0.001 & \textbf{144} & 0.003 & \textbf{144} & 0.003 & \textbf{144} & 0.160 & \textbf{144} & 0.018\\
c03 & 500 & 625 & 83 & 754 & 769 & 0.020 & 761 & 0.315 & 761 & 0.052 & -- & -- & \textbf{754} & 0.386\\
c04 & 500 & 625 & 125 & 1079 & 1105 & 0.015 & 1088 & 0.902 & 1100 & 0.119 & -- & -- & \textbf{1079} & 0.604\\
c05 & 500 & 625 & 250 & 1579 & 1585 & 0.036 & 1581 & 3.505 & 1581 & 0.705 & -- & -- & \textbf{1579} & 1.115\\
c06 & 500 & 1000 & 5 & 55 & 60 & 0.001 & \textbf{55} & 0.003 & \textbf{55} & 0.001 & \textbf{55} & 0.005 & \textbf{55} & 0.223\\
c07 & 500 & 1000 & 10 & 102 & 115 & 0.003 & \textbf{102} & 0.011 & 107 & 0.003 & \textbf{102} & 0.167 & \textbf{102} & 0.256\\
c08 & 500 & 1000 & 83 & 509 & 526 & 0.022 & 512 & 0.953 & 514 & 0.086 & -- & -- & \textbf{509} & 0.901\\
c09 & 500 & 1000 & 125 & 707 & 721 & 0.026 & 714 & 0.990 & 713 & 0.106 & -- & -- & \textbf{707} & 4.717\\
c10 & 500 & 1000 & 250 & 1093 & 1106 & 0.039 & 1095 & 4.338 & 1098 & 0.833 & -- & -- & \textbf{1093} & 3.277\\
c11 & 500 & 2500 & 5 & 32 & 37 & 0.002 & 34 & 0.004 & 33 & 0.002 & \textbf{32} & 0.005 & \textbf{32} & 0.506\\
c12 & 500 & 2500 & 10 & 46 & 49 & 0.003 & \textbf{46} & 0.028 & 48 & 0.003 & \textbf{46} & 0.159 & \textbf{46} & 0.419\\
c13 & 500 & 2500 & 83 & 258 & 268 & 0.029 & 260 & 1.049 & 264 & 0.041 & -- & -- & \textbf{258} & 9.081\\
c14 & 500 & 2500 & 125 & 323 & 337 & 0.038 & 328 & 2.465 & 331 & 0.175 & -- & -- & \textbf{323} & 3.840\\
c15 & 500 & 2500 & 250 & 556 & 561 & 0.094 & 557 & 4.952 & 559 & 0.574 & -- & -- & \textbf{556} & 4.400\\
c16 & 500 & 12500 & 5 & 11 & 12 & 0.002 & 12 & 0.005 & 12 & 0.003 & \textbf{11} & 0.007 & \textbf{11} & 0.451\\
c17 & 500 & 12500 & 10 & 18 & 20 & 0.005 & 19 & 0.017 & 19 & 0.005 & \textbf{18} & 0.238 & \textbf{18} & 0.465\\
c18 & 500 & 12500 & 83 & 113 & 128 & 0.037 & 120 & 0.827 & 119 & 0.050 & -- & -- & 117\textsuperscript{t} & 300.0\\
c19 & 500 & 12500 & 125 & 146 & 154 & 0.026 & 155 & 2.756 & 148 & 0.084 & -- & -- & \textbf{146} & 7.929\\
c20 & 500 & 12500 & 250 & 267 & 268 & 0.121 & 268 & 10.9 & \textbf{267} & 0.230 & -- & -- & \textbf{267} & 4.765\\
\bottomrule\end{longtable}\endgroup
\subsection*{Set P6E}
\begingroup\scriptsize\setlength{\tabcolsep}{3.2pt}
\begin{longtable}{@{}lrrrrrrrrrrrrrr@{}}\toprule
Inst. & \(n\) & \(m\) & \(k\) & OPT & \multicolumn{2}{c}{G1} & \multicolumn{2}{c}{G2} & \multicolumn{2}{c}{G3} & \multicolumn{2}{c}{DP} & \multicolumn{2}{c}{B\&B}\\
 & & & & & cost & s & cost & s & cost & s & cost & s & cost & s\\\midrule\endhead
p619 & 100 & 180 & 5 & 7485 & \textbf{7485} & 0.000 & \textbf{7485} & 0.000 & \textbf{7485} & 0.000 & \textbf{7485} & 0.001 & \textbf{7485} & 0.007\\
p620 & 100 & 180 & 5 & 8746 & 8853 & 0.000 & \textbf{8746} & 0.000 & \textbf{8746} & 0.000 & \textbf{8746} & 0.001 & \textbf{8746} & 0.006\\
p621 & 100 & 180 & 5 & 8688 & 8791 & 0.000 & \textbf{8688} & 0.000 & \textbf{8688} & 0.000 & \textbf{8688} & 0.001 & \textbf{8688} & 0.007\\
p622 & 100 & 180 & 10 & 15972 & 17403 & 0.000 & \textbf{15972} & 0.001 & \textbf{15972} & 0.000 & \textbf{15972} & 0.022 & \textbf{15972} & 0.009\\
p623 & 100 & 180 & 10 & 19496 & 20667 & 0.000 & \textbf{19496} & 0.002 & \textbf{19496} & 0.000 & \textbf{19496} & 0.023 & \textbf{19496} & 0.011\\
p624 & 100 & 180 & 20 & 20246 & \textbf{20246} & 0.001 & \textbf{20246} & 0.005 & \textbf{20246} & 0.001 & \textbf{20246} & 150.0 & \textbf{20246} & 0.045\\
p625 & 100 & 180 & 20 & 23078 & 24193 & 0.001 & \textbf{23078} & 0.012 & 23098 & 0.001 & \textbf{23078} & 155.1 & \textbf{23078} & 0.039\\
p626 & 100 & 180 & 20 & 22346 & 23000 & 0.003 & \textbf{22346} & 0.014 & \textbf{22346} & 0.002 & \textbf{22346} & 140.8 & \textbf{22346} & 0.043\\
p627 & 100 & 180 & 50 & 40647 & 40748 & 0.003 & \textbf{40647} & 0.073 & \textbf{40647} & 0.001 & -- & -- & \textbf{40647} & 0.088\\
p628 & 100 & 180 & 50 & 40008 & 40607 & 0.001 & \textbf{40008} & 0.041 & \textbf{40008} & 0.002 & -- & -- & \textbf{40008} & 0.140\\
p629 & 100 & 180 & 50 & 43287 & 43746 & 0.001 & 43335 & 0.035 & \textbf{43287} & 0.002 & -- & -- & \textbf{43287} & 0.107\\
p630 & 200 & 370 & 10 & 26125 & 26761 & 0.001 & 26416 & 0.002 & \textbf{26125} & 0.001 & \textbf{26125} & 0.047 & \textbf{26125} & 0.033\\
p631 & 200 & 370 & 20 & 39067 & 42103 & 0.001 & 40429 & 0.011 & 39565 & 0.002 & -- & -- & \textbf{39067} & 2.407\\
p632 & 200 & 370 & 40 & 56217 & 57776 & 0.003 & 56562 & 0.052 & \textbf{56217} & 0.003 & -- & -- & \textbf{56217} & 13.1\\
p633 & 200 & 370 & 100 & 86268 & 86877 & 0.007 & 86297 & 0.336 & \textbf{86268} & 0.012 & -- & -- & \textbf{86268} & 0.523\\
\bottomrule\end{longtable}\endgroup
\subsection*{Set Adversarial}
\begingroup\scriptsize\setlength{\tabcolsep}{3.2pt}
\begin{longtable}{@{}lrrrrrrrrrrrrrr@{}}\toprule
Inst. & \(n\) & \(m\) & \(k\) & OPT & \multicolumn{2}{c}{G1} & \multicolumn{2}{c}{G2} & \multicolumn{2}{c}{G3} & \multicolumn{2}{c}{DP} & \multicolumn{2}{c}{B\&B}\\
 & & & & & cost & s & cost & s & cost & s & cost & s & cost & s\\\midrule\endhead
adv\_g1 & 7 & 8 & 4 & 30 & 42 & 0.000 & \textbf{30} & 0.000 & \textbf{30} & 0.000 & \textbf{30} & 0.000 & \textbf{30} & 0.000\\
adv\_g2 & 8 & 10 & 5 & 49 & \textbf{49} & 0.000 & 68 & 0.000 & \textbf{49} & 0.000 & \textbf{49} & 0.000 & \textbf{49} & 0.000\\
adv\_g3 & 9 & 9 & 5 & 48 & 58 & 0.000 & \textbf{48} & 0.000 & 58 & 0.000 & \textbf{48} & 0.000 & \textbf{48} & 0.000\\
star\_l03 & 4 & 6 & 3 & 30 & 38 & 0.000 & 38 & 0.000 & \textbf{30} & 0.000 & \textbf{30} & 0.000 & \textbf{30} & 0.000\\
star\_l04 & 5 & 10 & 4 & 40 & 57 & 0.000 & 57 & 0.000 & \textbf{40} & 0.000 & \textbf{40} & 0.000 & \textbf{40} & 0.000\\
star\_l05 & 6 & 15 & 5 & 50 & 76 & 0.000 & 76 & 0.000 & \textbf{50} & 0.000 & \textbf{50} & 0.000 & \textbf{50} & 0.000\\
star\_l06 & 7 & 21 & 6 & 60 & 95 & 0.000 & 95 & 0.000 & \textbf{60} & 0.000 & \textbf{60} & 0.000 & \textbf{60} & 0.000\\
star\_l07 & 8 & 28 & 7 & 70 & 114 & 0.000 & 114 & 0.000 & \textbf{70} & 0.000 & \textbf{70} & 0.000 & \textbf{70} & 0.000\\
star\_l08 & 9 & 36 & 8 & 80 & 133 & 0.000 & 133 & 0.000 & \textbf{80} & 0.000 & \textbf{80} & 0.000 & \textbf{80} & 0.000\\
star\_l09 & 10 & 45 & 9 & 90 & 152 & 0.000 & 152 & 0.001 & \textbf{90} & 0.000 & \textbf{90} & 0.001 & \textbf{90} & 0.000\\
star\_l10 & 11 & 55 & 10 & 100 & 171 & 0.000 & 171 & 0.000 & \textbf{100} & 0.000 & \textbf{100} & 0.001 & \textbf{100} & 0.001\\
star\_l11 & 12 & 66 & 11 & 110 & 190 & 0.000 & 190 & 0.000 & \textbf{110} & 0.000 & \textbf{110} & 0.005 & \textbf{110} & 0.001\\
star\_l12 & 13 & 78 & 12 & 120 & 209 & 0.000 & 209 & 0.000 & \textbf{120} & 0.000 & \textbf{120} & 0.007 & \textbf{120} & 0.001\\
\bottomrule\end{longtable}\endgroup
\subsection*{Set Grid}
\begingroup\scriptsize\setlength{\tabcolsep}{3.2pt}
\begin{longtable}{@{}lrrrrrrrrrrrrrr@{}}\toprule
Inst. & \(n\) & \(m\) & \(k\) & OPT & \multicolumn{2}{c}{G1} & \multicolumn{2}{c}{G2} & \multicolumn{2}{c}{G3} & \multicolumn{2}{c}{DP} & \multicolumn{2}{c}{B\&B}\\
 & & & & & cost & s & cost & s & cost & s & cost & s & cost & s\\\midrule\endhead
grid\_k04\_s0 & 100 & 180 & 4 & 89 & 94 & 0.000 & 90 & 0.000 & \textbf{89} & 0.000 & \textbf{89} & 0.000 & \textbf{89} & 0.023\\
grid\_k04\_s1 & 100 & 180 & 4 & 33 & \textbf{33} & 0.000 & \textbf{33} & 0.000 & \textbf{33} & 0.000 & \textbf{33} & 0.001 & \textbf{33} & 0.014\\
grid\_k04\_s2 & 100 & 180 & 4 & 34 & \textbf{34} & 0.000 & \textbf{34} & 0.000 & \textbf{34} & 0.000 & \textbf{34} & 0.000 & \textbf{34} & 0.010\\
grid\_k06\_s0 & 100 & 180 & 6 & 96 & 103 & 0.000 & \textbf{96} & 0.001 & \textbf{96} & 0.000 & \textbf{96} & 0.001 & \textbf{96} & 0.005\\
grid\_k06\_s1 & 100 & 180 & 6 & 77 & \textbf{77} & 0.000 & \textbf{77} & 0.001 & \textbf{77} & 0.000 & \textbf{77} & 0.001 & \textbf{77} & 0.013\\
grid\_k06\_s2 & 100 & 180 & 6 & 61 & \textbf{61} & 0.000 & 62 & 0.000 & \textbf{61} & 0.000 & \textbf{61} & 0.001 & \textbf{61} & 0.014\\
grid\_k08\_s0 & 100 & 180 & 8 & 117 & 118 & 0.000 & 118 & 0.000 & \textbf{117} & 0.000 & \textbf{117} & 0.002 & \textbf{117} & 0.005\\
grid\_k08\_s1 & 100 & 180 & 8 & 94 & 98 & 0.000 & \textbf{94} & 0.000 & 96 & 0.000 & \textbf{94} & 0.002 & \textbf{94} & 0.013\\
grid\_k08\_s2 & 100 & 180 & 8 & 85 & 93 & 0.000 & 86 & 0.001 & \textbf{85} & 0.000 & \textbf{85} & 0.005 & \textbf{85} & 0.005\\
grid\_k10\_s0 & 100 & 180 & 10 & 107 & 116 & 0.000 & \textbf{107} & 0.002 & 108 & 0.000 & \textbf{107} & 0.022 & \textbf{107} & 0.016\\
grid\_k10\_s1 & 100 & 180 & 10 & 113 & 115 & 0.000 & \textbf{113} & 0.002 & 114 & 0.000 & \textbf{113} & 0.022 & \textbf{113} & 0.017\\
grid\_k10\_s2 & 100 & 180 & 10 & 94 & \textbf{94} & 0.000 & \textbf{94} & 0.001 & \textbf{94} & 0.000 & \textbf{94} & 0.022 & \textbf{94} & 0.010\\
grid\_k12\_s0 & 100 & 180 & 12 & 98 & 101 & 0.000 & \textbf{98} & 0.002 & \textbf{98} & 0.000 & \textbf{98} & 0.095 & \textbf{98} & 0.015\\
grid\_k12\_s1 & 100 & 180 & 12 & 106 & 113 & 0.001 & \textbf{106} & 0.002 & 107 & 0.001 & \textbf{106} & 0.097 & \textbf{106} & 0.014\\
grid\_k12\_s2 & 100 & 180 & 12 & 121 & 122 & 0.000 & 122 & 0.002 & \textbf{121} & 0.001 & \textbf{121} & 0.097 & \textbf{121} & 0.020\\
grid\_k14\_s0 & 100 & 180 & 14 & 126 & 130 & 0.001 & 129 & 0.003 & 129 & 0.001 & \textbf{126} & 0.493 & \textbf{126} & 0.017\\
grid\_k14\_s1 & 100 & 180 & 14 & 147 & 156 & 0.001 & \textbf{147} & 0.003 & 151 & 0.001 & \textbf{147} & 0.462 & \textbf{147} & 0.014\\
grid\_k14\_s2 & 100 & 180 & 14 & 111 & 118 & 0.001 & 112 & 0.002 & 113 & 0.001 & \textbf{111} & 0.440 & \textbf{111} & 0.014\\
grid\_k16\_s0 & 100 & 180 & 16 & 121 & 125 & 0.001 & \textbf{121} & 0.004 & 123 & 0.001 & \textbf{121} & 1.363 & \textbf{121} & 0.008\\
grid\_k16\_s1 & 100 & 180 & 16 & 123 & 124 & 0.001 & \textbf{123} & 0.004 & \textbf{123} & 0.001 & \textbf{123} & 3.380 & \textbf{123} & 0.020\\
grid\_k16\_s2 & 100 & 180 & 16 & 109 & \textbf{109} & 0.001 & \textbf{109} & 0.003 & \textbf{109} & 0.001 & \textbf{109} & 3.531 & \textbf{109} & 0.013\\
grid\_k18\_s0 & 100 & 180 & 18 & 112 & 115 & 0.001 & 115 & 0.005 & \textbf{112} & 0.003 & \textbf{112} & 16.1 & \textbf{112} & 0.022\\
grid\_k18\_s1 & 100 & 180 & 18 & 119 & 121 & 0.001 & 121 & 0.004 & \textbf{119} & 0.001 & \textbf{119} & 14.3 & \textbf{119} & 0.021\\
grid\_k18\_s2 & 100 & 180 & 18 & 151 & 154 & 0.001 & 153 & 0.005 & 152 & 0.001 & \textbf{151} & 16.9 & \textbf{151} & 0.040\\
grid\_k20\_s0 & 100 & 180 & 20 & 166 & 169 & 0.001 & \textbf{166} & 0.012 & 168 & 0.001 & \textbf{166} & 124.7 & \textbf{166} & 0.020\\
grid\_k20\_s1 & 100 & 180 & 20 & 139 & 140 & 0.001 & \textbf{139} & 0.005 & \textbf{139} & 0.001 & \textbf{139} & 113.7 & \textbf{139} & 0.017\\
grid\_k20\_s2 & 100 & 180 & 20 & 151 & 153 & 0.003 & \textbf{151} & 0.003 & \textbf{151} & 0.001 & \textbf{151} & 107.7 & \textbf{151} & 0.024\\
\bottomrule\end{longtable}\endgroup
\subsection*{Set Random}
\begingroup\scriptsize\setlength{\tabcolsep}{3.2pt}
\begin{longtable}{@{}lrrrrrrrrrrrrrr@{}}\toprule
Inst. & \(n\) & \(m\) & \(k\) & OPT & \multicolumn{2}{c}{G1} & \multicolumn{2}{c}{G2} & \multicolumn{2}{c}{G3} & \multicolumn{2}{c}{DP} & \multicolumn{2}{c}{B\&B}\\
 & & & & & cost & s & cost & s & cost & s & cost & s & cost & s\\\midrule\endhead
rand\_n0100 & 100 & 200 & 10 & 689 & 699 & 0.000 & 697 & 0.002 & 691 & 0.001 & \textbf{689} & 0.024 & \textbf{689} & 0.011\\
rand\_n0200 & 200 & 400 & 10 & 827 & 945 & 0.001 & \textbf{827} & 0.004 & 864 & 0.001 & \textbf{827} & 0.058 & \textbf{827} & 0.033\\
rand\_n0400 & 400 & 800 & 10 & 1226 & 1370 & 0.002 & 1244 & 0.011 & 1291 & 0.002 & \textbf{1226} & 0.177 & \textbf{1226} & 0.176\\
rand\_n0800 & 800 & 1600 & 10 & 2005 & 2194 & 0.003 & 2046 & 0.018 & 2046 & 0.003 & \textbf{2005} & 0.246 & \textbf{2005} & 0.670\\
rand\_n1600 & 1600 & 3200 & 10 & 2652 & 2909 & 0.007 & 2691 & 0.034 & 2706 & 0.007 & \textbf{2652} & 0.525 & \textbf{2652} & 2.086\\
rand\_n3200 & 3200 & 6400 & 10 & 4556 & 4606 & 0.014 & 4557 & 0.065 & 4574 & 0.014 & \textbf{4556} & 1.116 & \textbf{4556} & 8.735\\
\bottomrule\end{longtable}\endgroup
\subsection*{Set I080 (100 instances, grouped)}
\begin{table}[H]\centering\small\begin{tabular}{@{}rrrrrrrrrrrr@{}}\toprule
\(k\) & \# & \multicolumn{2}{c}{G1} & \multicolumn{2}{c}{G2} & \multicolumn{2}{c}{G3} & \multicolumn{2}{c}{DP} & \multicolumn{2}{c}{B\&B}\\
 & & gap \% & opt & gap \% & opt & gap \% & opt & solved & max s & solved & max s\\\midrule
6 & 25 & 14.83 & 0 & 7.76 & 5 & 7.53 & 2 & 25 & 0.01 & 25 & 0.03\\
8 & 25 & 18.28 & 1 & 8.42 & 5 & 6.44 & 1 & 25 & 0.01 & 25 & 0.04\\
16 & 25 & 19.66 & 0 & 12.83 & 1 & 9.27 & 0 & 25 & 4.40 & 25 & 0.29\\
20 & 25 & 21.27 & 0 & 13.93 & 2 & 9.32 & 0 & 25 & 159.31 & 25 & 0.53\\
\bottomrule\end{tabular}\end{table}


